% sec4_6.tex — Section 4.6: Common Coefficients

\section{Common Coefficients}
\label{sec:4.6}

In many applications, particularly in panel data contexts, you deal with a special
case of the multiple-equation model where the number of regressors is the same
across equations with the same coefficients. Such a model is a special case of the
multiple-equation model of Section~4.1. This section shows how to apply GMM
to multiple equations while imposing the common coefficient restriction. It will be
shown at the end that the seemingly restrictive model actually includes as a special
case the multiple-equation model \textit{without} the common coefficient restriction.

\subsection*{The Model with Common Coefficients}

With common coefficients, the multiple-equation system becomes

\begin{quote}
\textbf{Assumption 4.1$'$ (linearity):}\quad
\textit{There are $M$ linear equations to be estimated:}
\begin{equation}
y_{im} = \mathbf{z}_{im}' \boldsymbol{\delta} + \varepsilon_{im}
\quad (m = 1, 2, \ldots, M;\; i = 1, 2, \ldots, n),
\label{eq:4.6.1}
\end{equation}
\textit{where $n$ is the sample size, $\mathbf{z}_{im}$ is an $L$-dimensional vector of
regressors, $\boldsymbol{\delta}$ is an $L$-dimensional coefficient vector common to all
equations, and $\varepsilon_{im}$ is an unobservable error term of the $m$-th equation.}
\end{quote}

Of the other assumptions of the multiple-equation model, Assumptions~4.2--4.6,
only Assumption~4.4 (identification) needs to be modified to reflect the common
coefficient restriction. The multiple-equation version of
$\mathbf{g}(\mathbf{w}_i;\, \tilde{\boldsymbol{\delta}})$ is now
\begin{equation}
\mathbf{g}(\mathbf{w}_i;\, \tilde{\boldsymbol{\delta}})
\equiv \begin{bmatrix}
\mathbf{x}_{i1} \cdot (y_{i1} - \mathbf{z}_{i1}' \tilde{\boldsymbol{\delta}}) \\
\vdots \\
\mathbf{x}_{iM} \cdot (y_{iM} - \mathbf{z}_{iM}' \tilde{\boldsymbol{\delta}})
\end{bmatrix},
\label{eq:4.6.2}
\end{equation}
so that $\E[\mathbf{g}(\mathbf{w}_i;\, \tilde{\boldsymbol{\delta}})]$ becomes
\begin{equation}
\E[\mathbf{g}(\mathbf{w}_i;\, \tilde{\boldsymbol{\delta}})]
= \begin{bmatrix}
\E(\mathbf{x}_{i1} \cdot y_{i1}) \\ \vdots \\ \E(\mathbf{x}_{iM} \cdot y_{iM})
\end{bmatrix}
- \begin{bmatrix}
\E(\mathbf{x}_{i1} \cdot \mathbf{z}_{i1}') \tilde{\boldsymbol{\delta}} \\
\vdots \\
\E(\mathbf{x}_{iM} \cdot \mathbf{z}_{iM}') \tilde{\boldsymbol{\delta}}
\end{bmatrix}
= \underset{(\sum_{m=1}^{M} K_m \times 1)}{\boldsymbol{\sigma}_{\mathbf{xy}}}
- \underset{(\sum_{m=1}^{M} K_m \times L)}{\boldsymbol{\Sigma}_{\mathbf{xz}}}
\underset{(L \times 1)}{\tilde{\boldsymbol{\delta}}},
\label{eq:4.6.3}
\end{equation}
where
\begin{equation}
\boldsymbol{\sigma}_{\mathbf{xy}}
\equiv \begin{bmatrix}
\E(\mathbf{x}_{i1} \cdot y_{i1}) \\ \vdots \\ \E(\mathbf{x}_{iM} \cdot y_{iM})
\end{bmatrix}, \quad
\boldsymbol{\Sigma}_{\mathbf{xz}}
\equiv \begin{bmatrix}
\E(\mathbf{x}_{i1} \mathbf{z}_{i1}') \\ \vdots \\ \E(\mathbf{x}_{iM} \mathbf{z}_{iM}')
\end{bmatrix}.
\label{eq:4.6.4}
\end{equation}
So the implication of common coefficients is that $\boldsymbol{\Sigma}_{\mathbf{xz}}$
is now a stacked, not a block diagonal, matrix. With this change, the system of
equations determining $\tilde{\boldsymbol{\delta}}$ is again~\eqref{eq:4.1.10}, so the
identification condition is

\begin{quote}
\textbf{Assumption 4.4$'$ (identification with common coefficients):}\quad
\textit{The $\sum_{m=1}^{M} K_m \times L$ matrix $\boldsymbol{\Sigma}_{\mathbf{xz}}$
defined in~\eqref{eq:4.6.4} is of full column rank.}
\end{quote}

This condition is weaker than Assumption~4.4, which requires that each equation
of the system be identified. Indeed, a sufficient condition for identification is that
$\E(\mathbf{x}_{im} \mathbf{z}_{im}')$ be of full column rank for \textit{some}
$m$.\footnote{Proof: If $\E(\mathbf{x}_{im} \mathbf{z}_{im}')$ is of full column rank,
its rank is $L$. Since the rank of a matrix is also equal to the number of linearly
independent rows, $\E(\mathbf{x}_{im} \mathbf{z}_{im}')$ has $L$ linearly independent
rows. Thus, $\boldsymbol{\Sigma}_{\mathbf{xz}}$ has at least $L$ linearly independent rows.}
This is thanks to the \textit{a priori} restriction that the coefficients are the same
across equations. It is possible that the system is identified even if none of the
equations are identified individually.

\subsection*{The GMM Estimator}

It is left to you to verify that $\mathbf{g}_n(\tilde{\boldsymbol{\delta}})$
can be written as $\mathbf{s}_{\mathbf{xy}} - \mathbf{S}_{\mathbf{xz}} \tilde{\boldsymbol{\delta}}$ with
\begin{equation}
\mathbf{s}_{\mathbf{xy}}
\equiv \begin{bmatrix}
\dfrac{1}{n} \displaystyle\sum_{i=1}^{n} \mathbf{x}_{i1} \cdot y_{i1} \\[6pt]
\vdots \\[4pt]
\dfrac{1}{n} \displaystyle\sum_{i=1}^{n} \mathbf{x}_{iM} \cdot y_{iM}
\end{bmatrix}, \quad
\mathbf{S}_{\mathbf{xz}}
\equiv \begin{bmatrix}
\dfrac{1}{n} \displaystyle\sum_{i=1}^{n} \mathbf{x}_{i1} \mathbf{z}_{i1}' \\[6pt]
\vdots \\[4pt]
\dfrac{1}{n} \displaystyle\sum_{i=1}^{n} \mathbf{x}_{iM} \mathbf{z}_{iM}'
\end{bmatrix},
\label{eq:4.6.5}
\end{equation}
so that the GMM estimator with a $\sum_m K_m \times \sum_m K_m$ weighting matrix
$\hat{\mathbf{W}}$ is~\eqref{eq:4.2.3} with the sampling error given
by~\eqref{eq:4.2.4}. Provided that $\boldsymbol{\Sigma}_{\mathbf{xz}}$
and $\mathbf{S}_{\mathbf{xz}}$ are as just redefined and $\boldsymbol{\delta}$
is interpreted as the common coefficient, large-sample theory for the
multiple-equation model with common coefficients is the same as that with separate
coefficients developed in Section~4.3, with Assumption~4.1$'$ replacing
Assumption~4.1, Assumption~4.4$'$ replacing Assumption~4.4, and the residual
$\hat{e}_i$ redefined to be $y_{im} - \mathbf{z}_{im}' \hat{\boldsymbol{\delta}}$
(not $y_{im} - \mathbf{z}_{im}' \hat{\boldsymbol{\delta}}_m$) for some consistent
estimate of the $L$-dimensional common coefficient vector $\boldsymbol{\delta}$.

In order to relate this GMM estimator to popular estimators in the literature,
it is necessary to write out the GMM formula in full. Substituting~\eqref{eq:4.6.5}
into~\eqref{eq:4.2.3}, we obtain
\begin{align}
\hat{\boldsymbol{\delta}}(\hat{\mathbf{W}})
&= \left(
\begin{bmatrix}
\dfrac{1}{n} \displaystyle\sum \mathbf{z}_{i1} \mathbf{x}_{i1}' & \cdots &
\dfrac{1}{n} \displaystyle\sum \mathbf{z}_{iM} \mathbf{x}_{iM}'
\end{bmatrix}
\begin{bmatrix}
\hat{\mathbf{W}}_{11} & \cdots & \hat{\mathbf{W}}_{1M} \\
\vdots & & \vdots \\
\hat{\mathbf{W}}_{M1} & \cdots & \hat{\mathbf{W}}_{MM}
\end{bmatrix}
\begin{bmatrix}
\dfrac{1}{n} \displaystyle\sum \mathbf{x}_{i1} \mathbf{z}_{i1}' \\[6pt]
\vdots \\[4pt]
\dfrac{1}{n} \displaystyle\sum \mathbf{x}_{iM} \mathbf{z}_{iM}'
\end{bmatrix}
\right)^{-1} \notag \\[6pt]
&\quad \times
\begin{bmatrix}
\dfrac{1}{n} \displaystyle\sum \mathbf{z}_{i1} \mathbf{x}_{i1}' & \cdots &
\dfrac{1}{n} \displaystyle\sum \mathbf{z}_{iM} \mathbf{x}_{iM}'
\end{bmatrix}
\begin{bmatrix}
\hat{\mathbf{W}}_{11} & \cdots & \hat{\mathbf{W}}_{1M} \\
\vdots & & \vdots \\
\hat{\mathbf{W}}_{M1} & \cdots & \hat{\mathbf{W}}_{MM}
\end{bmatrix}
\begin{bmatrix}
\dfrac{1}{n} \displaystyle\sum \mathbf{x}_{i1} \cdot y_{i1} \\[6pt]
\vdots \\[4pt]
\dfrac{1}{n} \displaystyle\sum \mathbf{x}_{iM} \cdot y_{iM}
\end{bmatrix} \notag \\[6pt]
&= \left[\sum_{m=1}^{M} \sum_{h=1}^{M}
\Bigl\{\Bigl(\frac{1}{n} \sum_{i=1}^{n} \mathbf{z}_{im} \mathbf{x}_{im}'\Bigr)
\hat{\mathbf{W}}_{mh}
\Bigl(\frac{1}{n} \sum_{i=1}^{n} \mathbf{x}_{ih} \mathbf{z}_{ih}'\Bigr)
\Bigr\}\right]^{-1} \notag \\
&\quad \times
\left[\sum_{m=1}^{M} \sum_{h=1}^{M}
\Bigl\{\Bigl(\frac{1}{n} \sum_{i=1}^{n} \mathbf{z}_{im} \mathbf{x}_{im}'\Bigr)
\hat{\mathbf{W}}_{mh}
\Bigl(\frac{1}{n} \sum_{i=1}^{n} \mathbf{x}_{ih} \cdot y_{ih}\Bigr)
\Bigr\}\right],
\label{eq:4.6.6}
\end{align}
where $\hat{\mathbf{W}}_{mh}$ is the $(m,h)$ block of $\hat{\mathbf{W}}$.
(For the second equality, use formula (A.8) of Appendix~A.)
The efficient GMM estimator obtains if the $\hat{\mathbf{W}}$ in this expression
is replaced by the inverse of $\widehat{\mathbf{S}}$ defined in~\eqref{eq:4.3.2}.

\subsection*{Imposing Conditional Homoskedasticity}

The efficient GMM estimator obtains when we set in~\eqref{eq:4.6.6}
$\hat{\mathbf{W}} = \widehat{\mathbf{S}}^{-1}$ where $\widehat{\mathbf{S}}$
is a consistent estimator of $\mathbf{S}$. As before, we can impose conditional
homoskedasticity to reduce the efficient GMM estimator to popular estimators.
The FIVE estimator obtains if $\widehat{\mathbf{S}}$ is given by~\eqref{eq:4.5.3}.

If we further assume that the set of instruments is the same across equations,
then, as before, this $\widehat{\mathbf{S}}$ has the Kronecker product structure
\[
\widehat{\mathbf{S}} = \widehat{\boldsymbol{\Sigma}} \otimes
\Bigl(\frac{1}{n} \sum \mathbf{x}_i \mathbf{x}_i'\Bigr),
\]
so that $\hat{\mathbf{W}}_{mh}$ in~\eqref{eq:4.6.6} is given by~\eqref{eq:4.5.11},
resulting in what should be called the 3SLS estimator with common coefficients:
\begin{equation}
\left[\sum_{m=1}^{M} \sum_{h=1}^{M}
\Bigl\{\hat{\sigma}^{mh} \cdot
\Bigl(\frac{1}{n} \sum \mathbf{z}_{im} \mathbf{x}_i'\Bigr)
\Bigl(\frac{1}{n} \sum \mathbf{x}_i \mathbf{x}_i'\Bigr)^{-1}
\Bigl(\frac{1}{n} \sum \mathbf{x}_i \mathbf{z}_{ih}'\Bigr)
\Bigr\}\right]^{-1}
\left[\sum_{m=1}^{M} \sum_{h=1}^{M}
\Bigl\{\hat{\sigma}^{mh} \cdot
\Bigl(\frac{1}{n} \sum \mathbf{z}_{im} \mathbf{x}_i'\Bigr)
\Bigl(\frac{1}{n} \sum \mathbf{x}_i \mathbf{x}_i'\Bigr)^{-1}
\Bigl(\frac{1}{n} \sum \mathbf{x}_i \cdot y_{ih}\Bigr)
\Bigr\}\right],
\label{eq:4.6.7}
\end{equation}
where $\hat{\sigma}^{mh}$ is the $(m,h)$ element of
$\widehat{\boldsymbol{\Sigma}}^{-1}$.
If, in addition, the SUR condition~\eqref{eq:4.5.18} is assumed, then the
``disappearance of $\mathbf{x}$'' occurs in~\eqref{eq:4.6.7} and the efficient
GMM estimator becomes what is called (for historical reasons) the
\textbf{random-effects estimator}:
\begin{equation}
\hat{\boldsymbol{\delta}}_{\mathrm{RE}}
= \left[\sum_{m=1}^{M} \sum_{h=1}^{M}
\hat{\sigma}^{mh} \Bigl(\frac{1}{n} \sum_{i=1}^{n}
\mathbf{z}_{im} \mathbf{z}_{ih}'\Bigr)\right]^{-1}
\sum_{m=1}^{M} \sum_{h=1}^{M}
\hat{\sigma}^{mh} \Bigl(\frac{1}{n} \sum_{i=1}^{n}
\mathbf{z}_{im} \cdot y_{ih}\Bigr).
\label{eq:4.6.8}
\end{equation}
Its asymptotic variance is
\begin{equation}
\avar(\hat{\boldsymbol{\delta}}_{\mathrm{RE}})
= \left[\sum_{m=1}^{M} \sum_{h=1}^{M}
\sigma^{mh}\, \E(\mathbf{z}_{im} \mathbf{z}_{ih}')\right]^{-1}.
\label{eq:4.6.9}
\end{equation}
(To derive this, go back to the general formula~\eqref{eq:4.3.3} for the Avar of
efficient GMM. Set $\boldsymbol{\Sigma}_{\mathbf{xz}}$ as in~\eqref{eq:4.6.4},
$\mathbf{S} = \boldsymbol{\Sigma} \otimes \E(\mathbf{x}_i \mathbf{x}_i')$,
use formula (A.8) of Appendix~A to calculate the matrix product, and observe the
``disappearance of $\mathbf{x}$'' that~\eqref{eq:4.5.16} becomes~\eqref{eq:4.5.16p}
on page~280 under the SUR assumption~\eqref{eq:4.5.18}.)
It is consistently estimated by
\begin{equation}
\widehat{\avar}(\hat{\boldsymbol{\delta}}_{\mathrm{RE}})
= \left[\sum_{m=1}^{M} \sum_{h=1}^{M}
\hat{\sigma}^{mh} \Bigl(\frac{1}{n} \sum_{i=1}^{n}
\mathbf{z}_{im} \mathbf{z}_{ih}'\Bigr)\right]^{-1}.
\label{eq:4.6.10}
\end{equation}

We summarize the result about the random-effects estimator as

\begin{proposition}[large-sample properties of the random-effects estimator]
\label{prop:4.7}
\textit{Suppose Assumptions~4.1$'$, 4.2, 4.3, 4.4$'$, 4.5, and~4.7 hold with
$\mathbf{x}_i = \text{union of}\; (\mathbf{z}_{i1}, \ldots, \mathbf{z}_{iM})$.
Let\/ $\boldsymbol{\Sigma}$ be the $M \times M$ matrix whose $(m,h)$ element is
$\E(\varepsilon_{im}\, \varepsilon_{ih})$, and let\/ $\widehat{\boldsymbol{\Sigma}}$
be a consistent estimate of\/ $\boldsymbol{\Sigma}$. Then}
\begin{enumerate}[label=(\alph*)]
\item $\hat{\boldsymbol{\delta}}_{\mathrm{RE}}$, \textit{given by~\eqref{eq:4.6.8}
is consistent, asymptotically normal, and efficient, with\/
$\avar(\hat{\boldsymbol{\delta}}_{\mathrm{RE}})$ given by~\eqref{eq:4.6.9}.}

\item \textit{The estimated asymptotic variance~\eqref{eq:4.6.10} is consistent for\/
$\avar(\hat{\boldsymbol{\delta}}_{\mathrm{RE}})$.}

\item \textit{(Sargan's statistic)}
\[
J(\hat{\boldsymbol{\delta}}_{\mathrm{RE}},\, \widehat{\mathbf{S}}^{-1})
\equiv n \cdot \mathbf{g}_n(\hat{\boldsymbol{\delta}}_{\mathrm{RE}})'
\widehat{\mathbf{S}}^{-1}\, \mathbf{g}_n(\hat{\boldsymbol{\delta}}_{\mathrm{RE}})
\dto \chi^2(MK - L),
\]
\textit{where\/ $\widehat{\mathbf{S}} = \widehat{\boldsymbol{\Sigma}} \otimes
\bigl(\frac{1}{n} \sum_i \mathbf{x}_i \mathbf{x}_i'\bigr)$,
$K$ is the number of common instruments, and\/
$\mathbf{g}_n(\tilde{\boldsymbol{\delta}}) = \mathbf{s}_{\mathbf{xy}} -
\mathbf{S}_{\mathbf{xz}} \tilde{\boldsymbol{\delta}}$ with\/
$\mathbf{s}_{\mathbf{xy}}$ and\/ $\mathbf{S}_{\mathbf{xz}}$ given
by~\eqref{eq:4.6.5}.}
\end{enumerate}
\end{proposition}

Below we will rewrite the formulas~\eqref{eq:4.6.8}--\eqref{eq:4.6.10} in more
elegant forms.

\subsection*{Pooled OLS}

For the SUR of the previous section, we obtained $\widehat{\boldsymbol{\Sigma}}$
from the residuals from the equation-by-equation OLS. The consistency is all that
is required for $\widehat{\boldsymbol{\Sigma}}$, so the same procedure works here
just as well, but the finite sample distribution of
$\hat{\boldsymbol{\delta}}_{\mathrm{RE}}$ might be improved if we exploit the
\textit{a priori} restriction that the coefficients be the same across equations
in the estimation of $\widehat{\boldsymbol{\Sigma}}$.
So consider setting $\hat{\mathbf{W}}$ to
\[
\mathbf{I}_M \otimes \Bigl(\frac{1}{n} \sum \mathbf{x}_i \mathbf{x}_i'\Bigr)^{-1},
\]
rather than
\[
\widehat{\boldsymbol{\Sigma}}^{-1} \otimes
\Bigl(\frac{1}{n} \sum \mathbf{x}_i \mathbf{x}_i'\Bigr)^{-1},
\]
in the first-step GMM estimation for the purpose of an initial consistent
estimate of $\boldsymbol{\delta}$. The estimator is~\eqref{eq:4.6.8} with
$\hat{\sigma}^{mh} = 1$ for $m = h$ and $0$ for $m \neq h$,
which can be written as
\begin{equation}
\hat{\boldsymbol{\delta}}_{\text{pooled OLS}}
= \left[\sum_{m=1}^{M} \Bigl(\frac{1}{n} \sum_{i=1}^{n}
\mathbf{z}_{im} \mathbf{z}_{im}'\Bigr)\right]^{-1}
\sum_{m=1}^{M} \Bigl(\frac{1}{n} \sum_{i=1}^{n}
\mathbf{z}_{im} \cdot y_{im}\Bigr)
= \Bigl(\sum_{i=1}^{n} \sum_{m=1}^{M}
\mathbf{z}_{im} \mathbf{z}_{im}'\Bigr)^{-1}
\sum_{i=1}^{n} \sum_{m=1}^{M} \mathbf{z}_{im} \cdot y_{im},
\label{eq:4.6.11}
\end{equation}
which is simply the OLS estimator on the sample of size $Mn$ where observations
are pooled across equations. For this reason, the estimator is called the
\textbf{pooled OLS} estimator. The expression suggests that the orthogonality conditions
that are being exploited are
\begin{equation}
\E(\mathbf{z}_{i1} \cdot \varepsilon_{i1} + \cdots
+ \mathbf{z}_{iM} \cdot \varepsilon_{iM}) = \mathbf{0},
\label{eq:4.6.12}
\end{equation}
which does not involve the ``cross'' orthogonalities
$\E(\mathbf{z}_{im} \cdot \varepsilon_{ih}) = \mathbf{0}$ $(m \neq h)$.
It is left as an analytical exercise to show that the pooled OLS is the GMM estimator
exploiting~\eqref{eq:4.6.12}.

Because it is so easy to calculate, pooled OLS is a popular option for those
researchers who do not want to program estimators. The estimator is also robust
to the failure of the ``cross'' orthogonalities. But it is important to keep in mind
that the standard errors printed out by the OLS package, which do not take into
account the interequation cross moments ($\sigma_{mh}$), are biased.
For the pooled OLS estimator, which is a GMM estimator with a nonoptimal choice of
$\hat{\mathbf{W}}$, the correct formula for the asymptotic variance is (3.5.1)
of Proposition~3.1. Setting
$\mathbf{W} = \mathbf{I}_M \otimes [\E(\mathbf{x}_i \mathbf{x}_i')]^{-1}$,
$\mathbf{S} = \boldsymbol{\Sigma} \otimes \E(\mathbf{x}_i \mathbf{x}_i')$,
$\boldsymbol{\Sigma}_{\mathbf{xz}}$ as in~\eqref{eq:4.6.4}, and again observing the
disappearance of $\mathbf{x}$, we obtain
\begin{equation}
\avar(\hat{\boldsymbol{\delta}}_{\text{pooled OLS}})
= \left[\sum_{m=1}^{M} \E(\mathbf{z}_{im} \mathbf{z}_{im}')\right]^{-1}
\sum_{m=1}^{M} \sum_{h=1}^{M}
\sigma_{mh}\, \E(\mathbf{z}_{im} \mathbf{z}_{ih}')
\left[\sum_{m=1}^{M} \E(\mathbf{z}_{im} \mathbf{z}_{im}')\right]^{-1},
\label{eq:4.6.13}
\end{equation}
which is consistently estimated by
\begin{equation}
\Bigl(\sum_{m=1}^{M} \frac{1}{n} \sum_{i=1}^{n}
\mathbf{z}_{im} \mathbf{z}_{im}'\Bigr)^{-1}
\sum_{m=1}^{M} \sum_{h=1}^{M}
\hat{\sigma}_{mh} \frac{1}{n} \sum_{i=1}^{n}
\mathbf{z}_{im} \mathbf{z}_{ih}'
\Bigl(\sum_{m=1}^{M} \frac{1}{n} \sum_{i=1}^{n}
\mathbf{z}_{im} \mathbf{z}_{im}'\Bigr)^{-1}.
\label{eq:4.6.14}
\end{equation}
Since the pooled OLS estimator is consistent, the associated residual can be used to
calculate $\hat{\sigma}_{mh}$ $(m, h = 1, \ldots, M)$ in this expression.
The correct standard errors of pooled OLS are the square roots of $(1/n$ times)
the diagonal element of this matrix.

\subsection*{Beautifying the Formulas}

These formulas for the random-effects and pooled OLS estimators are quite straightforward
to derive but look complicated, with double and triple summations. They
can, however, be beautified---and at the same time more useful for programming---if
we introduce some new matrix notation. Let
\begin{equation}
\underset{(M \times 1)}{\mathbf{y}_i}
= \begin{bmatrix} y_{i1} \\ \vdots \\ y_{iM} \end{bmatrix}, \quad
\underset{(M \times L)}{\mathbf{Z}_i}
= \begin{bmatrix} \mathbf{z}_{i1}' \\ \vdots \\ \mathbf{z}_{iM}' \end{bmatrix}, \quad
\underset{(M \times 1)}{\boldsymbol{\varepsilon}_i}
= \begin{bmatrix} \varepsilon_{i1} \\ \vdots \\ \varepsilon_{iM} \end{bmatrix},
\label{eq:4.6.15}
\end{equation}
so that the $M$-equation system with common coefficients in Assumption~4.1$'$ can
be written compactly as
\begin{equation}
\mathbf{y}_i = \mathbf{Z}_i \boldsymbol{\delta} + \boldsymbol{\varepsilon}_i
\quad (i = 1, 2, \ldots, n).
\tag{4.6.1$'$}
\label{eq:4.6.1p}
\end{equation}

The beautification of the complicated formulas utilizes the following algebraic
results:
\begin{align}
\sum_{m=1}^{M} \mathbf{z}_{im} \mathbf{z}_{im}' &= \mathbf{Z}_i' \mathbf{Z}_i, \quad
\sum_{m=1}^{M} \mathbf{z}_{im} \cdot y_{im} = \mathbf{Z}_i' \mathbf{y}_i,
\label{eq:4.6.16a} \\
\sum_{m=1}^{M} \sum_{h=1}^{M} c_{mh} \cdot \mathbf{z}_{im} \cdot y_{ih}
&= \mathbf{Z}_i' \mathbf{C} \mathbf{y}_i, \quad
\sum_{m=1}^{M} \sum_{h=1}^{M} c_{mh} \cdot \mathbf{z}_{im} \mathbf{z}_{ih}'
= \mathbf{Z}_i' \mathbf{C} \mathbf{Z}_i
\label{eq:4.6.16b}
\end{align}
for any $M \times M$ matrix $\mathbf{C} = (c_{mh})$.
Now take the triple summation in~\eqref{eq:4.6.8}
\begin{align}
\sum_{m=1}^{M} \sum_{h=1}^{M} \hat{\sigma}^{mh} \cdot
\Bigl(\frac{1}{n} \sum_{i=1}^{n} \mathbf{z}_{im} \cdot y_{ih}\Bigr)
&= \frac{1}{n} \sum_{i=1}^{n}
\Bigl(\sum_{m=1}^{M} \sum_{h=1}^{M} \hat{\sigma}^{mh} \cdot
\mathbf{z}_{im} \cdot y_{ih}\Bigr)
\quad \text{(by changing order of summations)} \notag \\
&= \frac{1}{n} \sum_{i=1}^{n} \mathbf{Z}_i'
\widehat{\boldsymbol{\Sigma}}^{-1} \mathbf{y}_i
\quad \text{(by~\eqref{eq:4.6.16b} with
$\mathbf{C} = \widehat{\boldsymbol{\Sigma}}^{-1}$).}
\label{eq:4.6.17}
\end{align}

Similarly, by~\eqref{eq:4.6.16b} with
$\mathbf{C} = \widehat{\boldsymbol{\Sigma}}^{-1}$, the other triple summation
in~\eqref{eq:4.6.8} becomes
\begin{equation}
\sum_{m=1}^{M} \sum_{h=1}^{M} \hat{\sigma}^{mh}
\frac{1}{n} \sum \mathbf{z}_{im} \mathbf{z}_{ih}'
= \frac{1}{n} \sum_{i=1}^{n} \mathbf{Z}_i'
\widehat{\boldsymbol{\Sigma}}^{-1} \mathbf{Z}_i.
\label{eq:4.6.18}
\end{equation}

The double and triple summations in~\eqref{eq:4.6.9} and~\eqref{eq:4.6.10}
can be written similarly. We have thus rewritten the random-effects formulas as
\begin{align}
\hat{\boldsymbol{\delta}}_{\mathrm{RE}}
&= \Bigl(\frac{1}{n} \sum_{i=1}^{n} \mathbf{Z}_i'
\widehat{\boldsymbol{\Sigma}}^{-1} \mathbf{Z}_i\Bigr)^{-1}
\frac{1}{n} \sum_{i=1}^{n} \mathbf{Z}_i'
\widehat{\boldsymbol{\Sigma}}^{-1} \mathbf{y}_i,
\tag{4.6.8$'$} \label{eq:4.6.8p} \\
\avar(\hat{\boldsymbol{\delta}}_{\mathrm{RE}})
&= \bigl(\E(\mathbf{Z}_i' \boldsymbol{\Sigma}^{-1} \mathbf{Z}_i)\bigr)^{-1},
\tag{4.6.9$'$} \label{eq:4.6.9p} \\
\widehat{\avar}(\hat{\boldsymbol{\delta}}_{\mathrm{RE}})
&= \Bigl(\frac{1}{n} \sum_{i=1}^{n} \mathbf{Z}_i'
\widehat{\boldsymbol{\Sigma}}^{-1} \mathbf{Z}_i\Bigr)^{-1}.
\tag{4.6.10$'$} \label{eq:4.6.10p}
\end{align}

Using~\eqref{eq:4.6.16a}, the pooled OLS formulas can be written as
\begin{align}
\hat{\boldsymbol{\delta}}_{\text{pooled OLS}}
&= \Bigl(\frac{1}{n} \sum_{i=1}^{n} \mathbf{Z}_i' \mathbf{Z}_i\Bigr)^{-1}
\frac{1}{n} \sum_{i=1}^{n} \mathbf{Z}_i' \mathbf{y}_i,
\tag{4.6.11$'$} \label{eq:4.6.11p} \\
\avar(\hat{\boldsymbol{\delta}}_{\text{pooled OLS}})
&= [\E(\mathbf{Z}_i' \mathbf{Z}_i)]^{-1}\,
\E(\mathbf{Z}_i' \boldsymbol{\Sigma} \mathbf{Z}_i)\,
[\E(\mathbf{Z}_i' \mathbf{Z}_i)]^{-1},
\tag{4.6.13$'$} \label{eq:4.6.13p} \\
\widehat{\avar}(\hat{\boldsymbol{\delta}}_{\text{pooled OLS}})
&= \Bigl(\frac{1}{n} \sum \mathbf{Z}_i' \mathbf{Z}_i\Bigr)^{-1}
\Bigl(\frac{1}{n} \sum \mathbf{Z}_i'
\widehat{\boldsymbol{\Sigma}} \mathbf{Z}_i\Bigr)
\Bigl(\frac{1}{n} \sum \mathbf{Z}_i' \mathbf{Z}_i\Bigr)^{-1}.
\tag{4.6.14$'$} \label{eq:4.6.14p}
\end{align}

These formulas are not just pretty; as will be seen in the next section, they will be
useful for handling certain classes of panel data (called \textbf{unbalanced panel data}).

\subsection*{The Restriction That Isn't}

Although it appears that the model of this section, with the common coefficient
assumption, is a special case of the model of Section~4.1, the latter can be cast in
the format of the model of this section with a proper redefinition of regressors.
For example, consider Example~4.1. Instead of defining $\mathbf{z}_{i1}$ and
$\mathbf{z}_{i2}$ as in~\eqref{eq:4.1.2}, define them as
\[
\mathbf{z}_{i1} = \begin{bmatrix}
1 \\ S_i \\ IQ_i \\ EXPR_i \\ 0 \\ 0 \\ 0
\end{bmatrix}, \quad
\mathbf{z}_{i2} = \begin{bmatrix}
0 \\ 0 \\ 0 \\ 0 \\ 1 \\ S_i \\ IQ_i
\end{bmatrix} \quad \text{with} \quad
\boldsymbol{\delta} = \begin{bmatrix}
\phi_1 \\ \beta_1 \\ \gamma_1 \\ \pi \\ \phi_2 \\ \beta_2 \\ \gamma_2
\end{bmatrix}.
\]
Then the two-equation system fits the format of Assumption~4.1$'$.
(See Review Question~5 below for a demonstration of the same point more generally.)

Furthermore, a system in which only a subset of variables have common coefficients
can be written in the form of Assumption~4.1$'$. Consider Example~4.2.
Suppose we wish to assume \textit{a priori} that the coefficients of schooling, $IQ$, and
experience remain constant over time (so $\beta_1 = \beta_2 = \beta$,
$\gamma_1 = \gamma_2 = \gamma$, $\pi_1 = \pi_2 = \pi$) but the intercept does not.
The translation is accomplished by setting
\begin{equation}
\mathbf{z}_{i1} = \begin{bmatrix} 1 \\ 0 \\ S69_i \\ IQ_i \\ EXPR69_i \end{bmatrix}, \quad
\mathbf{z}_{i2} = \begin{bmatrix} 0 \\ 1 \\ S80_i \\ IQ_i \\ EXPR80_i \end{bmatrix}, \quad
\boldsymbol{\delta} = \begin{bmatrix} \phi_1 \\ \phi_2 \\ \beta \\ \gamma \\ \pi \end{bmatrix}.
\label{eq:4.6.19}
\end{equation}

Therefore, the common coefficient restriction is not restrictive. We could have
derived the GMM estimator for the model of Section~4.1 as a special case of the
GMM estimator with the common coefficient restriction. But the model without
the restriction is perhaps an easier introduction to multiple equations.
